\section{A prime-independent quotient and combined-precision packets}
\label{sec-jp-packets}

Let \(n>324\) be prime, \(B=n^4\), \(L=\log(8B)\), and
\[
 K=\mathbb Q(\zeta),\qquad \zeta^2-\zeta+1=0,\qquad
 w_k=3k+\zeta,\qquad \alpha_k=w_k/\bar w_k
 \quad(B\le k<2B).
\]
Here the indices are integers. Fix a subset \(I\) on which the actual
\(\alpha_k\) are multiplicatively independent, and write
\[
 w_k^n=A_k+C_k\zeta,\qquad
 T_k=\lvert A_kC_k(A_k+C_k)\rvert .
\]
The previously proved primitive-power arithmetic makes these three arms
nonzero and pairwise coprime. Put
\[
 G=\{z\in K^\times:z\bar z=1\},\qquad
 \mu_3=\{1,\zeta^2,\zeta^4\},\qquad a_k=[\alpha_k]\in G/\mu_3.
\]
This section keeps a single actual quotient while using distinct finite
reductions. It does not estimate the number or total weight of different
primes with different owners.

\begin{lemma}[Actual quotient and finite reductions]
\label{jp-quotient}
The elements \(a_k\), \(k\in I\), are multiplicatively independent.
For every prime \(q>8B\) and positive integer \(e\), the quotient of the
norm-one residue group modulo \(q^e\) by its embedded \(\mu_3\) has an
\(n\)-torsion subgroup of cardinality at most \(n\).
If \(q^e\mid T_k\), the actual element \(a_k\) reduces into this subgroup.
\end{lemma}
\begin{proof}
A relation \(\prod_k\alpha_k^{x_k}\in\mu_3\) becomes
\(\prod_k\alpha_k^{3x_k}=1\) after cubing. Independence of the actual
\(\alpha_k\) gives \(x_k=0\) for every integer exponent.

Reduction is applied only to ratios whose numerator and denominator are
\(q\)-units. Every marked root has this property, because a prime dividing
one primitive boundary arm does not divide the norm. At precision one
the norm-one group is cyclic of order \(q-\chi_q\), where
\(\chi_q=(-3/q)\). Its quotient by the three distinct roots of unity has
order \((q-\chi_q)/3\). At higher precision the reduction kernel is an
abelian \(q\)-group. Since \(q\nmid n\), taking \(n\)-th powers is invertible
on that kernel. Thus \(n\)-torsion injects into the precision-one quotient
and has cardinality at most \(n\).

The actual arm identities give
\(\alpha_k^n\in\{1,\zeta^2,\zeta^4\}\pmod{q^e}\) when \(q^e\mid T_k\).
Consequently the quotient class has \(n\)-th power one. This argument
works in both the split and inert residue algebras.
\end{proof}

\begin{theorem}[Combined-precision rigidity]
\label{jp-rigidity}
Let \(S\) be a nonempty finite set of distinct primes \(q>8B\).
Assign positive integer precisions \(e_q\), set
\[
 m=\prod_{q\in S}q^{e_q},
\]
and suppose \(J\subseteq I\) satisfies
\[
 q^{e_q}\mid T_k\qquad\hbox{for every }k\in J
                   \hbox{ and every }q\in S.
\]
For a positive integer \(\nu\), assume
\[
 \log m\ge 6\nu L.
\]
Then the integer exponent vectors \(x=(x_k)_{k\in J}\) with
\(\sum_k|x_k|\le\nu\) inject into the product of the finite
\(n\)-torsion groups in Lemma~\ref{jp-quotient}.
\end{theorem}
\begin{proof}
Represent the actual product \(\prod_k\alpha_k^{x_k}\) as
\(z_x/\bar z_x\), using \(w_k\) for positive occurrences and
\(\bar w_k\) for negative occurrences. The zero vector uses the empty
product \(1\). These nonzero Eisenstein integers satisfy
\[
 |z_x|\le(6B)^\nu
\]
and are units at every marked prime.

Suppose \(x,y\) have equal images in all the finite quotient groups.
For each \(q\in S\) there is a unique \(\tau_q\in\mu_3\) such that their
ratios differ by \(\tau_q\) modulo \(q^{e_q}\).
For each of the three global \(\tau\), choose an Eisenstein unit \(v_\tau\)
with \(v_\tau/\bar v_\tau=\tau\). Let \(D_\tau\) be the integer coordinate
determinant of \(z_x\) and \(v_\tau z_y\).
The conjugate cross difference and \(q>3\) give
\(q^{e_q}\mid D_{\tau_q}\). Pairwise coprimality of the marked prime
powers therefore gives the single integer divisibility
\[
 m\mid D_1D_{\zeta^2}D_{\zeta^4}.
\]
No common value of \(\tau_q\) is assumed.

The actual Gram identity for Eisenstein coordinates gives
\[
 |D_\tau|
 \le\frac{2}{\sqrt3}|z_x|\,|z_y|
 \le\frac{2}{\sqrt3}(6B)^{2\nu}
 <(8B)^{2\nu}.
\]
The last strict inequality holds for every positive integer \(\nu\),
and multiplying by a unit does not change the norm. Hence
\[
 |D_1D_{\zeta^2}D_{\zeta^4}|<(8B)^{6\nu}\le m.
\]
An integer divisible by \(m\) and of smaller absolute value must vanish.
Some \(D_\tau\) is zero, so the exact algebraic ratios differ by that
global element of \(\mu_3\). Independence in Lemma~\ref{jp-quotient}
forces \(x=y\).

All signs, unequal product lengths and the empty vector are included.
Equivalently one may partition the primes according to their three
phases; some phase has product modulus at least \(m^{1/3}\).
The integer-product argument also includes an empty phase class.
\end{proof}

\begin{corollary}[Exact product-target capacity]
\label{jp-capacity}
With the hypotheses of Theorem~\ref{jp-rigidity}, put
\(M=|J|\) and \(s=|S|\). Then
\[
 \sum_{j=0}^{\min(M,\nu)}
       2^j\binom{M}{j}\binom{\nu}{j}\le n^s.
\]
\end{corollary}
\begin{proof}
Choose the \(j\) nonzero coordinates and their signs. Their positive
absolute values have sum at most \(\nu\), giving \(\binom{\nu}{j}\)
choices by adjoining a slack coordinate. The \(j=0\) term is the zero
vector. This counts precisely the domain of the injection; the product
of the target cardinalities is at most \(n^s\).
\end{proof}

The theorem combines precision from several primes, but every root of
\(J\) must satisfy every marked prime-power condition. Its domain is a
common-hit rectangle in the actual incidence relation. The size of the
product target cannot be replaced by \(n\) merely because every factor
has exponent \(n\). For \(M=1\) the displayed bound is only
\(2\nu+1\le n^s\); no claim of its optimality on a singleton is needed.
If a separate premise supplied a common phase for every comparison at
all marked primes, a single determinant and \(2\nu L\) would suffice.
That additional premise is not used here.

In particular the result neither bounds singleton prime labels nor
controls their weighted depths. Repeated fixed-prime arguments do not
produce a common finite target of cardinality \(n\). The full positive
and signed far tails, roots outside the independent domain, and global
ABC coverage remain unresolved. This section records the complete
ordinary JP proof; it makes no additional Lean claim about residue
groups, their actual arithmetic construction, or these open interfaces.
